\section*{Chapitre \ref{ch:m2:central-banks-and-the-short-rate} --- Banques centrales et taux court}

\begin{solution}{exo:m2:central-banks-and-the-short-rate:1}
Le jeudi compte pour 1 jour, le vendredi pour 3, le lundi pour 1~: $N = 5$.
$(1 + 0.0390/360)(1 + 0.0388 \times 3/360)(1 + 0.0392/360) - 1$, multiplié par
$360/5$, donne 3.8926~\%. La moyenne pondérée par les jours vaut $(3.90 + 3 \times 3.88 +
3.92)/5 = 3.8920\%$~; la capitalisation ajoute 0.06 points de base sur cinq jours.
\end{solution}

\begin{solution}{exo:m2:central-banks-and-the-short-rate:2}
Corridor $2.90 - 2.50 = 40$ points de base~; le taux de refinancement se situe 15 points de base
au-dessus du taux de dépôt. À 2.440~\%, le taux de référence est inférieur de 6 points de base
au taux de dépôt, ce qu'une banque ayant accès à la facilité ne ferait jamais
accepter comme prêteuse. Mais le taux mesure ce que les banques \emph{paient} pour emprunter,
y compris auprès des fonds monétaires, des assureurs et d'autres institutions qui
ne peuvent pas déposer auprès de la banque centrale~; la borne ne lie que les prêteurs ayant accès à la facilité.
\end{solution}

\begin{solution}{exo:m2:central-banks-and-the-short-rate:3}
Banque centrale~: titres $-10$, \omterm{def:m2:central-banks-and-the-short-rate:reserves}{réserves} $-10$. Banque de l'assureur~: \omterm{def:m2:central-banks-and-the-short-rate:reserves}{réserves}
$-10$, dépôt de l'assureur $-10$. Assureur~: obligations $+10$, dépôt $-10$. \omterm{def:m2:central-banks-and-the-short-rate:reserves}{Réserves}
et les dépôts sont détruits ensemble~; c'est le mécanisme du resserrement quantitatif
par cessions.
\end{solution}

\begin{solution}{exo:m2:central-banks-and-the-short-rate:4}
Approximation~: $r^2(N-1)/(2B) = 0.04^2 \times 90/720 = 2.00$ points de base.
Exactement~: $[(1 + 0.04/360)^{91} - 1] \times 360/91 - 0.04 = 2.01$ points de base.
Sur un trimestre, le terme n'est pas négligeable pour un swap de 1~milliard de dollars~: 2
points de base sur un trimestre représentent environ 50\,000~dollars.
\end{solution}

\begin{solution}{exo:m2:central-banks-and-the-short-rate:5}
Les \omterm{def:m2:central-banks-and-the-short-rate:reserves}{réserves} baissent de 150~milliards de dollars
(la \cref{prop:m2:central-banks-and-the-short-rate:total}). Dans un corridor,
l'offre se déplace vers la gauche le long de la partie pentue de la courbe de demande, et le taux au jour le jour
monte vers le taux de prêt à moins que le desk n'ajoute 150~milliards de dollars,
par repo ou par achats, le matin même~; il doit prévoir les flux fiscaux pour ce
faire. Dans un \omterm{def:m2:central-banks-and-the-short-rate:corridor}{système de plancher}, si les \omterm{def:m2:central-banks-and-the-short-rate:reserves}{réserves} restent sur la partie plate, le taux ne bouge pas et
la banque centrale ne fait rien. La taille du coussin qu'elle conserve est le prix de
ne pas avoir à prévoir.
\end{solution}

\begin{solution}{exo:m2:central-banks-and-the-short-rate:6}
Taux global $3.80 + 0.26161 + 1.50 = 5.56161\%$. Intérêts $50\,000\,000 \times
0.0556161 \times 91/360 = \text{dollars}~702\,926$. Les 26.161 points de base correspondent
au risque de crédit bancaire et à la prime de terme que contenait le LIBOR et que le SOFR ne contient pas.
\end{solution}

\begin{solution}{exo:m2:central-banks-and-the-short-rate:7}
Blocage de deux jours~: 3.7421~\%, soit 0.27 points de base de moins que le taux simple de 3.7448~\%
(le fixing de fin de trimestre du 30 septembre, 3.95~\%, est remplacé par celui du 28 septembre, 3.87~\%, et
celui du 29 septembre par le même). Décalage d'observation de cinq jours~: 3.6836~\%, 6.12
points de base de moins, soit exactement le chiffre du lookback. Les deux utilisent les fixings du 25
août au 23 septembre~; elles ne diffèrent que par le poids porté par chaque fixing,
et ici les deux listes de poids forment le même multi-ensemble pour chaque taux~: le
poids de quatre jours du week-end de la fête du Travail tombe sur un fixing à 3.62~\% dans les deux cas
(le 28 août dans le cas du lookback, le 4 septembre dans celui du décalage), et le
poids de trois jours font de même. Avec seulement deux niveaux de taux dans la fenêtre, le
produit défini par la \cref{def:m2:central-banks-and-the-short-rate:arrears} est
identique. Une hausse pendant le long week-end romprait l'égalité.
\end{solution}

\begin{solution}{exo:m2:central-banks-and-the-short-rate:8}
En 2026, la Fed a fait fonctionner un \omterm{def:m2:central-banks-and-the-short-rate:corridor}{système de plancher} avec des \omterm{def:m2:central-banks-and-the-short-rate:reserves}{réserves} abondantes~: elle a
relevé son \omterm{def:m2:central-banks-and-the-short-rate:policy}{taux directeur} en remontant les taux qu'elle administre (rémunération des \omterm{def:m2:central-banks-and-the-short-rate:reserves}{réserves}, taux du reverse repo et
taux de la \omterm{def:m2:central-banks-and-the-short-rate:policy}{facilité permanente} de repo, taux d'escompte) et n'a pas eu à modifier la
quantité de \omterm{def:m2:central-banks-and-the-short-rate:reserves}{réserves} du tout~; ses achats de bons du Trésor pendant les mêmes semaines
\emph{ajoutaient} des \omterm{def:m2:central-banks-and-the-short-rate:reserves}{réserves}. Le communiqué décrit un \omterm{def:m2:central-banks-and-the-short-rate:corridor}{système de corridor} avec des \omterm{def:m2:central-banks-and-the-short-rate:reserves}{réserves} rares,
où le desk retire ou ajoute chaque jour des \omterm{def:m2:central-banks-and-the-short-rate:reserves}{réserves} pour déplacer le taux
le long de la partie pentue de la courbe de demande, à peu près comme la Fed opérait avant
2008.
\end{solution}

\begin{solution}{pb:m2:central-banks-and-the-short-rate:1}
\textbf{1.} $2 \times 10^9 \times (0.0225 \times 15 + 0.0250 \times 15)/360 =
\text{euros}~3\,958\,333$.
\textbf{2.} $2 \times 10^9 \times (0.0265 \times 15 + 0.0290 \times 15)/360 =
\text{euros}~4\,625\,000$.
\textbf{3.} 666\,667~euros, versés à la banque centrale, qui rémunère au taux de dépôt
et perçoit le taux de prêt sur les mêmes \omterm{def:m2:central-banks-and-the-short-rate:reserves}{réserves}.
\textbf{4.} Afin que les banques préfèrent échanger entre elles~: l'écart récompense un
marché interbancaire, qui diffuse de l'information sur les banques et fait le
travail d'allocation de la banque centrale à sa place, et il fait de la facilité de prêt un
filet de sécurité plutôt qu'une habitude.
\textbf{5.} Tous les taux du marché se situent à l'intérieur du corridor
(la \cref{prop:m2:central-banks-and-the-short-rate:bounds})~; déplacer les deux bornes
déplace tout l'intervalle, et le pouvoir de négociation relatif des banques reste inchangé.
\textbf{6.} Milieu~: 2.45~\% puis 2.70~\%. Intérêts 4\,291\,667~euros.
\textbf{7.} A gagne 333\,333~euros par rapport à la facilité de dépôt~; B économise
333\,333~euros par rapport à la facilité de prêt.
\textbf{8.} La rareté des \omterm{def:m2:central-banks-and-the-short-rate:reserves}{réserves}, et donc l'alternative de chaque partie~:
avec davantage de \omterm{def:m2:central-banks-and-the-short-rate:reserves}{réserves} excédentaires dans le système, les prêteurs se font concurrence et le taux baisse
vers le taux de dépôt~; avec moins de \omterm{def:m2:central-banks-and-the-short-rate:reserves}{réserves}, ce sont les emprunteurs qui se font concurrence et le taux monte. Il s'agit
de la courbe de demande de la \cref{fig:m2:central-banks-and-the-short-rate:corridor}.
\textbf{9.} 5 points de base sur 2~milliards d'euros pendant 30 jours font 83\,333~euros~: A
gagne 416\,667~euros, B économise 250\,000~euros. La prime est le prix du
crédit de B, que la facilité ne facture pas, car elle exige un collatéral.
\textbf{10.} Un indice de référence sans collatéral (\texteuro STR, SONIA) mesure ce que A
facture à B sans collatéral~; un indice garanti (SOFR) mesure le marché du repo,
où le collatéral élimine l'essentiel de la prime de crédit.
\textbf{11.} Au taux de dépôt~: avec 50~milliards d'euros ajoutés, B détient des \omterm{def:m2:central-banks-and-the-short-rate:reserves}{réserves}
propres, personne n'a besoin d'emprunter, et la seule alternative d'un prêteur est la
facilité de dépôt.
\textbf{12.} Rien~: B n'emprunte pas, ou ne paierait pas plus que ce que A gagne à
la facilité.
\textbf{13.} Elle relève le taux de dépôt (et le taux de prêt avec lui)~; elle n'a plus
besoin d'ajuster la quantité de \omterm{def:m2:central-banks-and-the-short-rate:reserves}{réserves}.
\textbf{14.} Pas pour le taux. Son bilan constitue désormais un
instrument distinct~: pour les rendements à plus long terme, ou simplement un coussin qui maintient les \omterm{def:m2:central-banks-and-the-short-rate:reserves}{réserves} sur
la partie plate de la courbe, quelles que soient les variations des facteurs autonomes.
\textbf{15.} $10^9 \times 0.0006 \times 30/360 = \text{euros}~50\,000$~: A gagne
l'écart entre le taux du fonds et celui de la facilité, sans risque, car il
peut utiliser la facilité et le fonds ne le peut pas.
\textbf{16.} Un bilan plus petit et moins d'intérêts versés aux banques sur
les \omterm{def:m2:central-banks-and-the-short-rate:reserves}{réserves}~; un marché interbancaire actif où les banques se surveillent mutuellement et
le taux porte de l'information.
\textbf{17.} Le contrôle du taux sans prévoir les facteurs autonomes chaque
jour~; la liberté d'utiliser le bilan (achats après une crise, ou
de liquidité) sans perdre la maîtrise du taux.
\textbf{18.} Le refinancement n'étant qu'à 15 points de base au-dessus du taux de dépôt,
les banques en manque de \omterm{def:m2:central-banks-and-the-short-rate:reserves}{réserves} empruntent à bon compte auprès de la banque centrale, si bien que les \omterm{def:m2:central-banks-and-the-short-rate:reserves}{réserves} sont
offertes à la demande et le taux de marché reste proche du taux de dépôt même
lorsque le bilan se contracte~: un plancher, alimenté par la demande.
\textbf{19.} Résultat à retenir~: \emph{gain tiré des échanges de \omterm{def:m2:central-banks-and-the-short-rate:reserves}{réserves}} sur la
période est $(i_L - i_D) \times \text{montant} \times \text{jours}/360 = 0.40\%
\times \text{euros}~2\,\text{milliards} \times 30/360 = \text{euros}~666\,667$,
partagé entre A et B selon le niveau où se fixe le taux interbancaire. Dans un plancher, il
est nul~: aucune banque n'est en manque, et le taux se place sur le taux de dépôt.
\textbf{20.} Les taux que la banque centrale administre, et le point où l'offre
de \omterm{def:m2:central-banks-and-the-short-rate:reserves}{réserves} qu'elle choisit rencontre la demande des banques.
\end{solution}

\begin{solution}{iq:m2:central-banks-and-the-short-rate:1}
Le LIBOR était l'estimation, par un panel de banques, du taux auquel elles pouvaient emprunter
sans garantie pour une échéance (d'un à douze mois), fixé chaque matin pour la période
à venir. Le SOFR est un taux au jour le jour fondé sur des transactions de repo garanties
par des bons du Trésor américain, composé sur la période et connu à son terme. Le LIBOR
contenait le risque de crédit des banques et une prime de terme~; le SOFR ne contient ni l'un ni l'autre. Un prêt qui
passe au SOFR ajoute un \omterm{def:m2:central-banks-and-the-short-rate:fallback}{spread de repli} (26.161 points de base pour le LIBOR en dollars à trois mois) ou une
marge plus élevée afin que le prêteur ne perde rien.
\iqlookfor{garanti ou non, échéance ou jour le jour, prospectif ou a posteriori,
et pourquoi le spread existe.}
\end{solution}

\begin{solution}{iq:m2:central-banks-and-the-short-rate:2}
Le taux de dépôt ne contraint que les établissements pouvant recourir à la facilité de dépôt.
Les établissements riches en liquidités qui ne le peuvent pas (fonds monétaires, assureurs, certaines
agences publiques) prêtent aux banques, lesquelles n'acceptent ces liquidités qu'en dessous du
taux de dépôt et les déposent à la banque centrale pour la différence. Un
taux de référence non garanti qui intègre de tels emprunts, comme le \texteuro STR, s'affiche
sous le plancher~; la facilité de reverse repo de la Fed, ouverte aux fonds monétaires, place un
second plancher sous cette fuite.
\iqlookfor{qui a accès à la facilité, et l'arbitrage qui en découle.}
\end{solution}

\begin{solution}{iq:m2:central-banks-and-the-short-rate:3}
Banque centrale~: obligations en hausse, \omterm{def:m2:central-banks-and-the-short-rate:reserves}{réserves} en hausse. La banque du vendeur~: \omterm{def:m2:central-banks-and-the-short-rate:reserves}{réserves} en hausse, le
dépôt du vendeur en hausse. Le vendeur~: obligations en baisse, dépôt en hausse. Les \omterm{def:m2:central-banks-and-the-short-rate:reserves}{réserves} ne sont pas
«~prêtées~»~: une banque qui prête crée un dépôt, et les \omterm{def:m2:central-banks-and-the-short-rate:reserves}{réserves} ne circulent
qu'entre banques~; leur total est fixé par la banque centrale. Ce que modifie l'\omterm{def:m2:central-banks-and-the-short-rate:qe}{assouplissement quantitatif}, c'est la
duration et la composition du portefeuille détenu par le secteur privé, donc les rendements à plus long terme,
et le volume de \omterm{def:m2:central-banks-and-the-short-rate:reserves}{réserves}, donc la place du taux au jour le jour sur la courbe de
demande.
\iqlookfor{les trois comptes en T, et le rejet du récit du multiplicateur
monétaire.}
\end{solution}

\begin{solution}{iq:m2:central-banks-and-the-short-rate:4}
Le \omterm{def:m2:central-banks-and-the-short-rate:arrears}{composé a posteriori} ne connaît le taux qu'après le dernier jour de la période,
trop tard pour émettre une facture et un paiement ce jour-là. Un lookback utilise
les fixings de $p$ jours ouvrés antérieurs, avec les pondérations en jours propres à la période
d'intérêt~; un décalage d'observation recule toute la période d'observation de $p$
jours et utilise \emph{ses} pondérations. Le décalage est ce que reproduit un indice composé publié
(rapport de deux valeurs d'indice), ce qui le rend plus simple à vérifier~; le lookback
suit le calendrier du prêt. Elles donnent des résultats différents lorsque les week-ends et les jours fériés
tombent différemment dans les deux périodes.
\iqlookfor{la raison opérationnelle, et un énoncé précis de la différence
de pondération.}
\end{solution}

\begin{solution}{iq:m2:central-banks-and-the-short-rate:5}
Regarder où se place le taux de marché au jour le jour par rapport aux taux des \omterm{def:m2:central-banks-and-the-short-rate:policy}{facilités permanentes}
et s'il réagit aux variations des \omterm{def:m2:central-banks-and-the-short-rate:reserves}{réserves}. Dans un plancher, il se situe au niveau du taux de
dépôt ou juste en dessous, et bouge à peine lorsque les \omterm{def:m2:central-banks-and-the-short-rate:reserves}{réserves} varient de
montants importants (jours d'échéance fiscale, compte du Trésor)~; dans un corridor, il se place à l'intérieur
du corridor et bondit avec les chocs sur les \omterm{def:m2:central-banks-and-the-short-rate:reserves}{réserves} et en fin de trimestre et de mois.
Le bilan de la banque centrale (\omterm{def:m2:central-banks-and-the-short-rate:reserves}{réserves} très supérieures au niveau requis) et ses
opérations quotidiennes (peu nombreuses, voire aucune) le confirment.
\iqlookfor{écart au taux de dépôt et sensibilité à l'offre de \omterm{def:m2:central-banks-and-the-short-rate:reserves}{réserves} comme
les deux tests.}
\end{solution}

\begin{solution}{iq:m2:central-banks-and-the-short-rate:6}
Les week-ends et jours fériés (pondérations de 3 ou 4 jours~; le bon calendrier de jours fériés,
qui, pour le SOFR, est le calendrier des titres d'État américains)~; un lookback qui
remonte dans la période précédente~; un décalage d'observation en fin de mois
où les deux périodes comptent un nombre différent de jours ouvrés~; les lockouts~;
fixings manquants ou republiés~; taux négatifs~; le premier jour de la période étant un
jour férié~; la base de calcul des jours selon la devise~; l'arrondi du taux composé
et du montant d'intérêt, ainsi que d'un indice publié le cas échéant~; une période d'un
jour~; les dates de paiement postérieures à la fin de la période~; la cohérence avec
l'indice composé publié, à son arrondi près.
\iqlookfor{une liste systématique organisée par calendrier, conventions, données et arrondis,
et des tests qui comparent à une source indépendante.}
\end{solution}
